\chapter{Specifications}

\section{Technical Specifications}

\textbf{System Requirements}

Minimum Requirements
\begin{itemize}
    \item \textbf{Android:} Version 8.0 (API Level 26) or higher
    \item \textbf{iOS:} Version 12.0 or higher
    \item \textbf{RAM:} Minimum 3GB for optimal machine learning and Bird Encyclopedia performance
    \item \textbf{Storage:} 2GB available storage for application and ML models
    \item \textbf{Hardware:} Camera and microphone required for complete functionality
\end{itemize}

Recommended Specifications
\begin{itemize}
    \item \textbf{Android:} Version 10.0 or higher
    \item \textbf{iOS:} Version 14.0 or higher
    \item \textbf{RAM:} 4GB or more
    \item \textbf{Storage:} 4GB available storage
    \item \textbf{Processor:} Modern ARM architecture (2020 or later)
\end{itemize}

\textbf{Machine Learning Model Specifications}

Image Classification
\begin{itemize}
    \item \textbf{Object Detection:} SSD MobileNet V1 (MLKit)
    \item \textbf{Bird Classification:} Custom MobileNetV2 Model
    \item \textbf{Input Format:} RGB images, variable dimensions
    \item \textbf{Output:} Classification probabilities with bounding box coordinates
\end{itemize}

Audio Classification
\begin{itemize}
    \item \textbf{Primary Model:} BirdNET\_GLOBAL\_6K\_V2.4\_Model\_FP32.tflite
    \item \textbf{Meta Model:} BirdNET\_GLOBAL\_6K\_V2.4\_MData\_Model\_FP16.tflite
    \item \textbf{Input Format:} 48kHz, mono, 3-second audio samples
    \item \textbf{Output:} 6,522 bird species classification probabilities
    \item \textbf{Model Size:} Approximately 25-30 MB combined
\end{itemize}

\section{Development Environment Setup}

\textbf{Prerequisites}
\begin{lstlisting}[caption=Package Versions, label=lst:package_versions]
{
  "expo": "~51.0.39",
  "react": "18.2.0",
  "react-native": "0.74.5",
  "@tensorflow/tfjs": "^4.22.0",
  "@tensorflow/tfjs-react-native": "^1.0.0",
  "react-native-vision-camera": "^4.6.4",
  "@infinitered/react-native-mlkit-object-detection": "^3.1.0",
  "@infinitered/react-native-mlkit-image-labeling": "^3.1.0",
  "expo-sqlite": "~15.0.0",
  "firebase": "^11.6.1",
  "react-native-fast-tflite": "^1.6.1"
}
\end{lstlisting}

\textbf{Development Commands}
\begin{lstlisting}[caption=Development Scripts, label=lst:dev_scripts]
# TypeScript type checking
npm run typecheck

# Start development server
npm start

# Run on Android device
npm run android:device

# Run on iOS simulator  
npm run ios

# Execute test suite
npm test

# Production build generation
npm run build
\end{lstlisting}

\section{Database Architecture and Schema Design}

\textbf{SQLite Database Schema Implementation}
\begin{lstlisting}[caption=Database Schema, label=lst:db_schema]
-- Bird observation records (main spotting table)
CREATE TABLE bird_spottings (
    id               INTEGER PRIMARY KEY AUTOINCREMENT,
    image_uri        TEXT,
    video_uri        TEXT,
    audio_uri        TEXT,
    text_note        TEXT,
    gps_lat          REAL,
    gps_lng          REAL,
    date             TEXT,
    bird_type        TEXT,
    image_prediction TEXT,
    audio_prediction TEXT,
    synced           INTEGER DEFAULT 0,
    latinBirDex      TEXT
);

-- BirDex species database (30,000+ species)
CREATE TABLE birddex (
    species_code           TEXT    PRIMARY KEY,
    english_name           TEXT    NOT NULL,
    scientific_name        TEXT    NOT NULL,
    category               TEXT,
    family                 TEXT,
    order_                 TEXT,
    sort_v2024             TEXT,
    clements_v2024b_change TEXT,
    text_for_website_v2024b TEXT,
    range                  TEXT,
    extinct                TEXT,
    extinct_year           TEXT,
    sort_v2023             TEXT,
    de_name                TEXT,
    es_name                TEXT,
    ukrainian_name         TEXT,
    ar_name                TEXT
);

-- Database metadata and versioning
CREATE TABLE metadata (
    key   TEXT PRIMARY KEY,
    value TEXT
);

-- Indexes for performance optimization
CREATE INDEX idx_bird_spottings_date ON bird_spottings(date);
CREATE INDEX idx_bird_spottings_synced ON bird_spottings(synced);
CREATE INDEX idx_birddex_english ON birddex(english_name COLLATE NOCASE);
CREATE INDEX idx_birddex_scientific ON birddex(scientific_name COLLATE NOCASE);
CREATE INDEX idx_birddex_category ON birddex(category);
CREATE INDEX idx_birddex_family ON birddex(family);
\end{lstlisting}

\section{Cloud Architecture and API Specifications}

\textbf{Firebase Firestore Data Collections}
\begin{lstlisting}[caption=Firestore Collections, label=lst:firestore_collections]
// User profile management
users/{userId} {
    email: string,
    displayName: string,
    createdAt: timestamp,
    settings: {
        language: string,
        theme: string,
        notifications: boolean
    }
}

// Bird observation records
observations/{observationId} {
    userId: string,
    speciesName: string,
    commonName: string,
    scientificName: string,
    confidence: number,
    timestamp: timestamp,
    location: geopoint,
    imageUrl?: string,
    audioUrl?: string,
    notes?: string,
    metadata: {
        appVersion: string,
        deviceInfo: string
    }
}

// Media file management
media/{mediaId} {
    observationId: string,
    type: 'image' | 'audio',
    fileName: string,
    size: number,
    uploadedAt: timestamp,
    downloadUrl: string
}
\end{lstlisting}

\section{Performance Evaluation and System Benchmarks}

\textbf{Machine Learning Pipeline Performance Analysis}
\begin{table}[H]
\centering
\begin{tabular}{|l|l|l|l|}
\hline
\textbf{Operation} & \textbf{Low-End Device} & \textbf{Mid-Range Device} & \textbf{High-End Device} \\
\hline
Photo capture & 800ms & 500ms & 300ms \\
Object detection & 1200ms & 800ms & 400ms \\
Image classification & 2000ms & 1200ms & 600ms \\
Audio recording & 3000ms & 3000ms & 3000ms \\
Audio classification & 1500ms & 1000ms & 500ms \\
\textbf{Complete cycle} & \textbf{8.5s} & \textbf{6.5s} & \textbf{4.8s} \\
\hline
\end{tabular}
\caption{Machine Learning Pipeline Performance by Device Classification}
\label{tab:ml_performance}
\end{table}

\textbf{Memory Utilization Analysis}
\begin{table}[H]
\centering
\begin{tabular}{|l|l|}
\hline
\textbf{System Component} & \textbf{Memory Utilization} \\
\hline
Application base & ~50 MB \\
ML model framework & ~30 MB \\
BirDex database & ~25 MB \\
Image cache system & ~100 MB (configurable) \\
User-generated data & Variable \\
\textbf{Minimum total} & \textbf{~105 MB} \\
\hline
\end{tabular}
\caption{Memory Utilization by System Components}
\label{tab:memory_usage}
\end{table}

\section{System Configuration Framework}

\textbf{Application Configuration Parameters}
\begin{lstlisting}[caption=Configuration Options, label=lst:config_options]
export const Config = {
  // Machine learning pipeline configuration
  camera: {
    pipelineDelay: 2, // Seconds between processing cycles
    confidenceThreshold: 0.7, // Minimum confidence for data persistence
    maxDetections: 5, // Maximum concurrent detection objects
    imageQuality: 0.3, // Image quality for ML processing optimization
  },
  
  // Audio processing configuration
  audio: {
    recordingDuration: 3000, // Recording duration in milliseconds
    sampleRate: 48000, // Sample rate in Hz
    channels: 1, // Mono channel configuration
    format: 'wav',
  },
  
  // Database management configuration
  database: {
    pageSize: 20, // Elements per pagination page
    cacheSize: 100, // Cache size in MB for image storage
    syncInterval: 300000, // Synchronization interval: 5 minutes in ms
  },
  
  // User interface configuration
  ui: {
    theme: 'auto', // Theme options: 'light', 'dark', 'auto'
    language: 'auto', // Language code or 'auto' detection
    hapticFeedback: true,
    animations: true,
  }
};
\end{lstlisting}

\section{System Diagnostics and Troubleshooting Framework}

\textbf{Common System Issues and Resolution Procedures}

Machine Learning Pipeline Initialization Failure
\textbf{Symptoms:} No ML processing activity, UI displays "Offline" status

\textbf{Resolution procedures:}
\begin{enumerate}
    \item Verify camera permission grants
    \item Validate ML model loading integrity
    \item Execute application restart procedure
    \item Confirm device storage availability (minimum 1GB free space)
\end{enumerate}

Audio Classification System Malfunction
\textbf{Symptoms:} Absence of audio predictions in HUD display

\textbf{Resolution procedures:}
\begin{enumerate}
    \item Verify microphone permission configuration
    \item Test hardware microphone accessibility
    \item Validate BirdNET model loading status
    \item Confirm audio format compatibility requirements
\end{enumerate}

System Performance Degradation
\textbf{Symptoms:} Reduced processing speed, elevated memory utilization

\textbf{Optimization procedures:}
\begin{enumerate}
    \item Increase \texttt{Config.camera.pipelineDelay} parameter
    \item Reduce image quality settings for ML processing
    \item Decrease cache allocation size
    \item Terminate background application processes
\end{enumerate}

\textbf{Diagnostic Command Interface}
\begin{lstlisting}[caption=Debug Commands, label=lst:debug_commands]
// Pipeline logging monitoring
# Search for diagnostic patterns:
[UnifiedPipeline] State: capturing_image
[UnifiedPipeline] State: detecting_objects
[UnifiedPipeline] State: recording_audio

// Error validation procedures
[UnifiedPipeline] image error: [Error details]
[UnifiedPipeline] audio error: [Error details]

// Performance metric analysis
[Performance] Photo capture time: 500ms
[Performance] Detection count: 3
[Performance] Audio processing: 1000ms
\end{lstlisting}

\section{Production Deployment Framework}

\textbf{Android Application Package Build Process}
\begin{lstlisting}[caption=Android Build Process, label=lst:android_build]
# Dependency installation procedure
npm install --legacy-peer-deps

# Native project generation
npx expo prebuild --clean --platform android

# Local build process (when EAS unavailable)
npx expo run:android --device

# EAS build process (recommended approach)
eas build --platform android --profile production
\end{lstlisting}

\textbf{iOS Application Build Framework}
\begin{lstlisting}[caption=iOS Build Process, label=lst:ios_build]
# Dependency installation procedure
npm install --legacy-peer-deps

# Native project initialization
npx expo prebuild --clean --platform ios

# EAS build process
eas build --platform ios --profile production
\end{lstlisting}

\section{Licensing Framework and Attribution}

\textbf{Open Source License Compliance}
\begin{itemize}
    \item \textbf{LogChirpy:} AGPL-3.0 License
    \item \textbf{React Native:} MIT License
    \item \textbf{Expo:} MIT License
    \item \textbf{TensorFlow.js:} Apache 2.0 License
    \item \textbf{BirdNET Models:} Custom Academic License
    \item \textbf{MLKit:} Google Terms of Service
\end{itemize}

\textbf{Data Source Attribution}
\begin{itemize}
    \item \textbf{BirDex Database:} Compiled from Cornell Lab CSV 2024
    \item \textbf{Wikipedia:} taxonomic imagery
    \item \textbf{OpenStreetMap:} Cartographic data and geospatial services
\end{itemize}

\textbf{External Technical Contributions}
\begin{itemize}
    \item \textbf{BirdNET Team:} Original audio classification models (.h5 source files)
    \item \textbf{whoBIRD Project:} Original Kotlin implementation (ported to React Native by Marty)
    \item \textbf{Infinite Red:} MLKit React Native packages
    \item \textbf{Marc Wannenmacher:} React Native Vision Camera library development
\end{itemize}